\section{Introduction}

\subsection{Background and Motivation}
Virtual power plants (VPPs) aggregate distributed resources and coordinate them as a single operating entity. Their performance depends on accurate forecasts of load, photovoltaic (PV) generation, and wind power \cite{Ma2025,Cao2023,Qiu2024,Moreno2020,Wu2024}. In practice, these forecasts must be accurate and physically plausible. PV output should vanish at night, and wind output should respect basic turbine operating limits.

Early VPP forecasting methods were mainly statistical. More recently, recurrent models and Transformers have become standard tools for multivariate time-series prediction \cite{Zhou2021informer,Wu2021autoformer,Nie2023patchtst,Zeng2023}. Recent studies have also adapted these models to load forecasting and VPP scheduling \cite{Qiu2024,Ashraf2025,Du2025,Li2025MetaInformer}. EPSR has also reported load and wind forecasting methods that explicitly incorporate meteorological covariates and attention-based weighting of physical attributes \cite{Li2023PLF,Xiong2022Wind}. These models can reduce forecast error, but low error alone is not enough for VPP operation. Predictions should also remain consistent with basic operating conditions.

\subsection{Challenges of Data-Driven VPP Forecasting}
Recent advances in deep learning have improved time-series forecasting accuracy, yet their application to VPP prediction still involves several practical challenges \cite{Vle2025,Arabzadeh2025}:

\textbf{1) Output feasibility and activation consistency:} Purely data-driven models can achieve low error while still producing non-physical outputs, such as non-zero PV generation at night or negative renewable power. These violations may be small, but they reduce the value of the forecast for downstream operation.

\textbf{2) Limits of post-processing and static penalties:} Post-hoc clipping is useful at deployment, but it does not shape the model during training. Physics-informed penalties help, yet fixed penalties alone may not capture weather-driven activation changes in PV and wind generation \cite{Raissi2019,Li2026,Asinyo2025}. A practical model should therefore combine train-time guidance with flexible residual prediction.

\textbf{3) Need for interpretable weather-response behavior:} In deployment, it is useful to know whether the model reacts to weather in a physically reasonable way. Existing interpretability tools can reveal such patterns after training, but they do not necessarily encode them into the model itself \cite{Vle2025,Lovo2025,Arabzadeh2025}.

\subsection{Overview of This Work}
To address these issues, we propose \textbf{PhysFormer}, a physics-guided Transformer for physically consistent VPP forecasting. The model combines an explicit physical mapping stream, a bounded residual output mechanism, a physics-guided gating and alignment module (PGCC), and curriculum-based constraint training. In this work, PGCC is used for guidance and alignment rather than causal discovery.

The main contributions of this paper are as follows:

\begin{enumerate}
    \item We propose a \textbf{Bounded Physical Act-Residual (BPAR)} mechanism that combines theory-based baselines with residual corrections. Under the adopted normalization, it preserves differentiability while maintaining a structurally non-negative output form.

    \item We introduce a \textbf{Physics-Guided Gating and Alignment} module (PGCC) that injects weather-dependent priors into PV and wind activation. Its main role is to improve alignment and interpretability.

    \item We adopt a \textbf{Physics-Constrained Curriculum Training} strategy to stabilize optimization, and we evaluate the full framework on a real 3-year VPP dataset.
\end{enumerate}

Experiments show that PhysFormer eliminates boundary violations while remaining competitive with strong Transformer baselines in mean squared error. The results indicate a useful balance among accuracy, physical consistency, and engineering interpretability.


